% Faithful formal-publication transcription: TheoremA' and full main proof, pp16:17--16:21.
\section*{4.2 Search-to-decision reduction for Alternating Matrix Space Isometry}
\begin{authorclaim}{TheoremA'}
Given an oracle deciding Alternating Matrix Space Isometry, the task of finding an isometry between two alternating matrix spaces $\mathcal A,\mathcal B\in\Lambda(n,\mathbb F_q)$, if it exists, can be solved using at most $q^{O(n)}$ oracle queries each of size at most $O(n^2)$, and in time either $q^{O(n)}\cdot n!=q^{\widetilde O(n)}$, or $q^{O(n+m)}$.
\end{authorclaim}
\begin{proof}
We first present the gadget construction. Then based on this gadget, we present the search-to-decision reduction.
\paragraph{Gadget construction.} Let $\mathbf A=(A_1,\ldots,A_m)$ be an ordered linear basis of $\mathcal A$, and let $\mathtt A\in\mathrm T(n\times n\times m,\mathbb F_q)$ be the 3-way array constructed from $\mathbf A$, so we can write
\[\mathtt A=\begin{bmatrix}0&a_{1,2}&a_{1,3}&\cdots&a_{1,n}\\-a_{1,2}&0&a_{2,3}&\cdots&a_{2,n}\\-a_{1,3}&-a_{2,3}&0&\cdots&a_{3,n}\\\vdots&\vdots&\vdots&\ddots&\vdots\\-a_{1,n}&-a_{2,n}&-a_{3,n}&\cdots&0\end{bmatrix},\]
where $a_{i,j}\in\mathbb F^m$, $1\le i<j\le n$ thought of as a vector coming out of the page.
We first consider a 3-tensor $\widetilde{\mathtt A}_i$ constructed from $\mathtt A$, for any $1\le i\le n-1$, as $\widetilde{\mathtt A}_i=$
\begin{center}\includegraphics[width=.98\linewidth]{sources/search-gadget.pdf}\end{center}
where $e_{j,k}$ is the $(m+2n(j-1)+k)$th standard basis vector, and $f_{j,k}$ is the $(m+2ni+n(j-1)+k)$th standard basis vector. A pictorial description can be seen by combining Figure2 (for the $e_{j,k}$) and [32, Figure3] (for the $f_{j,k}$).
We claim the following.
\begin{authorclaim}{Claim10}
If there exist invertible matrices $P$ and $Q$ to satisfy $P^t\widetilde{\mathtt A}_iP=\widetilde{\mathtt B}_i^Q$, then $P$ must be in the form $\begin{bmatrix}P_{1,1}&0&0\\0&P_{2,2}&0\\P_{3,1}&P_{3,2}&P_{3,3}\end{bmatrix}$, where $P_{1,1}$ is a monomial matrix of size $i\times i$, $P_{2,2}$ is of size $(n-i)\times(n-i)$, and $P_{3,3}$ is of size $(2ni+n)\times(2ni+n)$.
Furthermore, there exist such $P$ and $Q$ if and only if $\mathcal A$ and $\mathcal B$ are isometric by a matrix of the form $\begin{bmatrix}P_{1,1}&0\\0&P_{2,2}\end{bmatrix}$ where $P_{1,1}$ is a monomial matrix of size $i\times i$.
\end{authorclaim}
\begin{proof}
This claim is immediate by combining the arguments for the FGS gadget [28] as used in [32], and the monomial-restricting gadget introduced in Section4.1. We only outline the argument and point out some subtle issues here.
First, observe that for the lateral slices of $\widetilde{\mathtt A}_i$:
\begin{itemize}
\item The first $i$ lateral slices have rank in $[2n,3n)$. Note that the rank is strictly less than $3n$ because some tube fibers (coming out of the page) are 0 in the upper-left $n\times n$ sub-array.
\item The next $n-i$ lateral slices have rank in $[n,2n)$.
\item The remaining $2ni+n$ lateral slices have rank in $[1,n)$ (since $i\ge1$.)
\end{itemize}
Because of the above, for $P$ and $Q$ to satisfy $P^t\widetilde{\mathtt A}_iP=\widetilde{\mathtt B}_i^Q$, $P$ must be in the required form.
It is the furthermore statement that requires certain care. The only if direction is straightforward: after observing that $P$ has to be of the above form, we can easily verify that $\begin{bmatrix}P_{1,1}&0\\0&P_{2,2}\end{bmatrix}$ is an isometry from $\mathcal A$ to $\mathcal B$. For the if direction, starting from $\begin{bmatrix}P_{1,1}&0\\0&P_{2,2}\end{bmatrix}$ and $Q_{1,1}\in\mathrm{GL}(m,\mathbb F)$, we need to design $P_{3,3}\in\mathrm{GL}(2ni+n,\mathbb F)$ and $Q_{2,2}\in\mathrm{GL}(2ni+n(n-i),\mathbb F)$ such that letting
\[P=\begin{bmatrix}P_{1,1}&0&0\\0&P_{2,2}&0\\0&0&P_{3,3}\end{bmatrix}\quad\text{and}\quad Q=\begin{bmatrix}Q_{1,1}&0\\0&Q_{2,2}\end{bmatrix},\]
we have $P^t\widetilde{\mathtt A}_iP=\widetilde{\mathtt B}_i^Q$. This can be achieved by combining the arguments for the only if directions in the proofs of Lemma8 and [32, Proposition3.3].
\end{proof}
\paragraph{The search-to-decision reduction.} Given these preparations, we now present the search-to-decision reduction for Alternating Matrix Space Isometry. Recall that this requires us to use the decision oracle $\mathcal O$ to compute an explicit isometry transformation $P\in\mathrm{GL}(n,q)$, if $\mathcal A$ and $\mathcal B$ are indeed isometric. Think of $P$ as sending the standard basis $(\vec e_1,\ldots,\vec e_n)$ to another basis $(v_1,\ldots,v_n)$, where $\vec e_i$ and $v_i$ are in $\mathbb F_q^n$.
In the first step, we guess $v_1$, the image of $\vec e_1$, and a complement subspace of $\langle v_1\rangle$, at the cost of $q^{O(n)}$. For each such guess, let $P_1$ be the matrix which sends $\vec e_1\mapsto v_1$ and sends $\langle\vec e_2,\ldots,\vec e_n\rangle$ to the chosen complementary subspace arbitrarily. We apply $P_1$ to $\mathtt A$, and still call the resulting 3-way array $\mathtt A$ in the following. Then construct $\widetilde{\mathtt A}_1$ and $\widetilde{\mathtt B}_1$, and feed these two instances to the oracle $\mathcal O$. Note that, since $P_{1,1}$ (using notation as above) must be monomial, any equivalence between $\widetilde{\mathtt A}_1$ and $\widetilde{\mathtt B}_1$ must preserve our choice of $v_1$ up to scale. Thus, clearly, if $\mathcal A$ and $\mathcal B$ are indeed isometric and we guess the correct image of $\vec e_1$, then the oracle $\mathcal O$ will return yes (and conversely).
In the second step, we guess $v_2$, the image of $\vec e_2$, and a complement subspace of $\langle v_2\rangle$ within $\langle\vec e_2,\ldots,\vec e_n\rangle$, at the cost of $q^{O(n)}$. Note here that the previous step guarantees that there is an isometry respecting the direct sum decomposition $\langle v_1\rangle\oplus\langle\vec e_2,\ldots,\vec e_n\rangle$, so we need only search for a complement of $v_2$ within $\langle\vec e_2,\ldots,\vec e_n\rangle$, and not a more general complement of $\langle v_1,v_2\rangle$ in all of $\mathbb F_q^n$. This is crucial for the runtime, as at the $n/2$ step, the latter strategy would result in searching through $q^{\Theta(n^2)}$ possibilities.
For each such guess, we apply the corresponding transformation to $\mathtt A$ (and again call the resulting 3-way array $\mathtt A$). Then construct $\widetilde{\mathtt A}_2$ and $\widetilde{\mathtt B}_2$, and feed these two instances to the oracle $\mathcal O$. Clearly, if $\mathcal A$ and $\mathcal B$ are indeed isometric and we guess the correct image of $\vec e_2$ (and $\vec e_1$ from the previous step), then the oracle $\mathcal O$ will return yes. However, there is a small caveat here, namely we may guess some image of $\vec e_2$, such that $\mathcal A$ and $\mathcal B$ are actually isometric by some matrix $P$ of the form $\begin{bmatrix}P_{1,1}&0\\0&P_{2,2}\end{bmatrix}$ where $P_{1,1}$ is a monomial matrix of size 2 (instead of the more desired diagonal matrix). But this is fine, as it still ensures $P_{1,1}$ to be monomial, which is the key property to keep. This means that our choices of $\{v_1,v_2\}$ is correct as a set up to scaling, so we proceed.
In general, in the $i$th step, we maintain the property that $\mathcal A$ and $\mathcal B$ are isometric by some $P=\begin{bmatrix}P_{1,1}&0\\0&P_{2,2}\end{bmatrix}$ where $P_{1,1}$ is a monomial matrix of size $(i-1)\times(i-1)$. We guess $v_i$, the image of $\vec e_i$ in $\langle\vec e_i,\ldots,\vec e_n\rangle$, and a complement subspace of $\langle v_i\rangle$ within $\langle\vec e_i,\ldots,\vec e_n\rangle$. This cost is $q^{O(n)}$. For each such guess, we apply the corresponding transformation to $\mathtt A$ (and call the resulting 3-way array $\mathtt A$). Then construct $\widetilde{\mathtt A}_i$ and $\widetilde{\mathtt B}_i$, and feed these two instances to the oracle $\mathcal O$. Once we guess correctly, we ensure that $\mathcal A$ and $\mathcal B$ are isometric by $P=\begin{bmatrix}P_{1,1}&0\\0&P_{2,2}\end{bmatrix}$ where $P_{1,1}$ is a monomial matrix of size $i\times i$.
So after the $(n-1)$th step, we know that $\mathcal A$ and $\mathcal B$ are isometric by a monomial transformation. As the number of all monomial transformations is $(q-1)^n\cdot n!\le q^n\cdot2^{n\log n}=q^{\widetilde O(n)}$, we can enumerate all monomial transformations and check correspondingly. This gives an algorithm in time $q^{\widetilde O(n)}$. By resorting to Proposition11 which solves Alternating Matrix Space Monomial Isometry in time $q^{O(n+m)}$, we have an algorithm in time $q^{O(n+m)}$.
Note that all the instances we feed into the oracle $\mathcal O$ are of size $O(n^2)$. This concludes the proof.
\end{proof}
\section*{4.3 A simply-exponential algorithm for monomial isometry of alternating matrix spaces}
We now state the algorithm for monomial isometry used in TheoremA'.
\begin{authorclaim}{Proposition11}
Let $\mathcal A,\mathcal B\le\Lambda(n,q)$ be $m$-dimensional. Then there exists a $q^{O(n+m)}$-time algorithm that decides whether $\mathcal A$ and $\mathcal B$ are monomially isometric, and if so, computes an explicit monomial isometry.
\end{authorclaim}
\begin{proof}
Let $\mathcal A,\mathcal B\le\Lambda(n,q)$ be two $m$-dimensional alternating matrix spaces. Clearly, by incurring a multiplicative factor of $q^n$, we can reduce to the problem of testing whether $\mathcal A$ and $\mathcal B$ are permutationally isometric, i.e. whether there exists a permutation matrix $T\in\mathrm{GL}(n,q)$, such that $T^t\mathcal A T=\mathcal B$. We will solve this problem in time $2^{O(n)}\cdot q^{O(m)}$. This would give an algorithm with total running time $q^n\cdot2^{O(n)}\cdot q^{O(m)}=q^{O(n+m)}$. The basic idea of the algorithm comes from Luks's dynamic programming technique for Hypergraph Isomorphism [53].
\paragraph{Reducing to a generalized linear code equivalence problem.} Suppose $\mathcal A=\langle A_1,\ldots,A_m\rangle$, and $\mathcal B=\langle B_1,\ldots,B_m\rangle$. Let $\mathtt A$ and $\mathtt B$ be the $n\times n\times m$ 3-way arrays formed by the given bases of $\mathcal A$ and $\mathcal B$. For $S\subseteq[n]$ of size $s$, let $(A_i)_S$ be the submatrix of $A_i$ with row and column indices in $S$. Then let $\mathtt A_S$ be the $s\times s\times m$ 3-way array formed by $((A_1)_S,\ldots,(A_m)_S)$. Similarly we can define $\mathtt B_S$ for $S\subseteq[n]$.
For each $S\subseteq[n]$ of size $s$, let $\operatorname{Iso}(\mathtt A_{[s]},\mathtt B_S)$ be the coset in $S_n\times\mathrm{GL}(m,q)$, such that $(A,B)\in S_n\times\mathrm{GL}(m,q)$ if and only if the natural action of $(A,B)$ sends $\mathtt A_{[s]}$ to $\mathtt B_S$. Since all the matrices are alternating, their diagonal entries are zero, and thus $\mathtt A_{\{i\}}$ and $\mathtt B_{\{i\}}$ are both the $1\times1\times m$ zero vector for any $i$. It follows that if $s=1$ and $S=\{i\}$, $\operatorname{Iso}(\mathtt A_{[1]},\mathtt B_S)=G\times\mathrm{GL}(m,q)$, where $G$ is the coset of $S_n$ consisting of permutations sending 1 to $i$.
Suppose we have computed $\operatorname{Iso}(\mathtt A_{[s]},\mathtt B_S)$ for all $s<t$. Fix $T\subseteq[n]$, $|T|=t$, and let us compute $\operatorname{Iso}(\mathtt A_{[t]},\mathtt B_T)$. For any $(A,B)\in\operatorname{Iso}(\mathtt A_{[t]},\mathtt B_T)$, $A$ sends $[t-1]$ to some $T'\subseteq T$ of size $t-1$. So in this case, $(A,B)\in\operatorname{Iso}(\mathtt A_{[t-1]},\mathtt B_{T'})$, which has been computed. Let $T\setminus T'=\{t'\}$. On the other hand, for $(A,B)\in\operatorname{Iso}(\mathtt A_{[t-1]},\mathtt B_{T'})$ to be in $\operatorname{Iso}(\mathtt A_t,\mathtt B_T)$, $(A,B)$ needs to send the $t$th horizontal slice of $\mathtt A_{[t]}$ to the $t'$th horizontal slice of $\mathtt B_T$.
We first identify $T'$ with $[t-1]$. We then note that every horizontal slice of $\mathtt A_{[t]}$ has a row of zeros. So the problem now becomes: given two $(t-1)\times m$ matrices $P$ and $Q$ over $\mathbb F_q$, decide whether $P$ and $Q$ are the same under $G\le S_{t-1}\times\mathrm{GL}(m,q)$. (Note that $G=\operatorname{Iso}(\mathtt A_{[t-1]},\mathtt B_{T'})$ from above.) Clearly, this is a generalization of the Linear Code Equivalence problem. Furthermore, if we could solve this problem in time $2^{O(n)}\cdot q^{O(m)}$, we would have achieved our original goal.
\paragraph{Solving the generalized linear code equivalence problem.} We solve the above problem again by a dynamic programming scheme as follows. For $R\subseteq[t-1]$ of size $r$, $P_R$ denotes the $r\times m$ submatrix of $P$ with row indices from $R$. Let $\operatorname{Iso}'(P_{[r]},Q_R)$ be the coset in $S_{t-1}\times\mathrm{GL}(m,q)$, such that $(C,D)\in\operatorname{Iso}'(P_{[r]},Q_R)$ if and only if the natural action of $(C,D)$ sends $P_{[r]}$ to $Q_R$. If $r=0$, then $\operatorname{Iso}'(P_\emptyset,Q_\emptyset)=G$ where $G\le S_{t-1}\times\mathrm{GL}(m,q)$ is given as an input.
Suppose we have computed $\operatorname{Iso}'(P_{[r]},Q_R)$ for any $r<u$. Fix $U\subseteq[t-1]$, $|U|=u$, and let us compute $\operatorname{Iso}'(P_{[u]},Q_U)$. For any $(C,D)\in\operatorname{Iso}'(P_{[u]},Q_U)$, $C$ sends $[u-1]$ to some $U'\subseteq U$ of size $u-1$. So in this case, $(A,B)\in\operatorname{Iso}(P_{[u-1]},Q_U)$, which has been computed. Let $U\setminus U'=\{u'\}$. On the other hand, for $(C,D)\in\operatorname{Iso}(P_{[u-1]},Q_{U'})$ to be in $\operatorname{Iso}(P_{[u]},Q_U)$, $D$ needs to send the $u$th row of $P_{[u]}$ to the $u'$th row of $Q_U$. This subcoset of $\operatorname{Iso}(P_{[u-1]},Q_{U'})$ can be computed in time $q^{O(m)}$, by treating $\mathrm{GL}(m,q)$ as a permutation group on $\mathbb F_q^m$. We then take a union over size-$(u-1)$ subsets $U'$ to obtain a generating set for $\operatorname{Iso}(P_{[u]},Q_U)$. If necessary, we can reduce the generating set size by applying the standard permutation group machinery, as our time bound is $2^{O(n)}\cdot q^{O(m)}$, which is quite generous.
\end{proof}
